\chapter{Zusammenfassung und Lessons Learned}
\label{cha:zusammenfassung}

\section{Projektergebnis}

Das Projekt InspectAir hat alle definierten Ziele erreicht. Es liegt ein funktionsfähiger Prototyp vor, der folgende Leistungsmerkmale aufweist:

\begin{itemize}[nosep]
\item Gleichzeitige Erfassung von fünf Luftqualitätsparametern (CO\textsubscript{2}, PM1.0/PM2.5/PM10, VOC, Temperatur, Luftfeuchtigkeit) mit einer Update-Rate von \SI{2}{\second}
\item Darstellung auf 4-Zoll-TFT-Display in fünf wählbaren UI-Designs mit vierstufiger Farbcodierung nach WHO/UBA-Grenzwerten
\item WebApp mit Echtzeit-Messwerten und 24-Stunden-Verlaufsdiagramm
\item Automatische Präsenzerkennung per 24\,GHz-Radar für Display-Dimming und Light Sleep
\item Modulare Software-Architektur mit 29~Quelldateien in dreischichtiger Struktur
\item 3D-gedrucktes PETG-Gehäuse mit optimiertem Airflow
\item Materialkosten von ca.\ 136\,\euro{} (wobei viele Kleinteile in Multipacks mit deutlich niedrigeren Stückkosten kamen)
\end{itemize}

Alle 28~definierten Testfälle wurden bestanden. Die Requirements Traceability weist jedem Requirement mindestens einen verifizierten Testfall zu. Das V-Modell wurde vollständig durchlaufen.

\section{Lessons Learned}

Die folgenden drei Bereiche haben sich als besonders lehrreich und prägend für das Projektergebnis erwiesen.

%=============================================================================
\subsection{Lesson 1: Display-Integration -- Die Suche nach den richtigen Libraries}
\label{sec:lesson_display}
%=============================================================================

\subsubsection{Ausgangssituation}

Die Integration des 4-Zoll ST7796S-Displays mit dem ESP32-S3 erwies sich als eine der größten technischen Herausforderungen des Projekts. Obwohl TFT-Displays auf ESP32-Basis in der Maker-Community weit verbreitet sind, ist die Kombination aus ESP32-S3, ST7796S-Controller und LVGL~9 weniger dokumentiert.

\subsubsection{Probleme und Lösungsweg}

\paragraph{Problem 1: TFT\_eSPI vs.\ LovyanGFX}
Zunächst wurde die weit verbreitete TFT\_eSPI-Library getestet. Diese zeigte jedoch Inkompatibilitäten mit dem ESP32-S3, insbesondere bei der SPI-DMA-Konfiguration. Nach mehreren Tagen erfolgloser Fehlersuche wurde auf LovyanGFX gewechselt -- eine japanische Library, die explizit für ESP32-S3 und ST7796S optimiert ist. Die Umstellung erforderte eine komplette Neuschreibung der Display-Initialisierung.

\begin{lstlisting}[language=C++,caption={LovyanGFX Konfiguration für ST7796S},label={lst:lgfx_config}]
class LGFX : public lgfx::LGFX_Device {
    lgfx::Panel_ST7796 _panel;
    lgfx::Bus_SPI _bus;
    lgfx::Light_PWM _light;
public:
    LGFX(void) {
        auto cfg = _bus.config();
        cfg.spi_host = SPI2_HOST;
        cfg.spi_mode = 0;
        cfg.freq_write = 40000000;  // 40 MHz
        cfg.pin_sclk = TFT_SCLK;
        cfg.pin_mosi = TFT_MOSI;
        cfg.pin_dc = TFT_DC;
        _bus.config(cfg);
        _panel.setBus(&_bus);
        // ... weitere Konfiguration
    }
};
\end{lstlisting}

\paragraph{Problem 2: LVGL~8 vs.\ LVGL~9}
Die meisten Online-Tutorials und Beispiele beziehen sich auf LVGL~8. LVGL~9 (erschienen 2024) hat jedoch eine grundlegend überarbeitete API. Funktionen wie \texttt{lv\_obj\_set\_style\_*} haben neue Signaturen, das Event-Handling wurde umgestellt, und die Display-Initialisierung erfordert einen neuen \texttt{lv\_display\_t}-Ansatz. Die Migration kostete weitere zwei Tage.

\paragraph{Problem 3: Font-Rendering mit deutschen Umlauten}
Standard-LVGL-Fonts unterstützen keine deutschen Umlaute. Für die Darstellung von "`Temperatur"', "`Luftfeuchtigkeit"' etc.\ mussten Custom Fonts mit dem LVGL Font Converter erstellt werden. Dabei stellte sich heraus, dass der Font Converter spezifische Unicode-Ranges erfordert:

\begin{lstlisting}[language=bash,numbers=none,caption={Font-Generierung mit Umlauten}]
lv_font_conv --font arial.ttf --size 20 \
  --range 0x20-0x7F,0xC4,0xD6,0xDC,0xE4,0xF6,0xFC,0xDF \
  --format lvgl --output font_arial_20.c
\end{lstlisting}

\subsubsection{Erkenntnisse}

\begin{enumerate}[nosep]
\item \textbf{Dokumentation kritisch prüfen:} Viele Tutorials sind veraltet. Immer die offizielle Dokumentation und GitHub-Issues der aktuellen Library-Version konsultieren.
\item \textbf{Hardware-spezifische Libraries bevorzugen:} Speziell optimierte Libraries (wie LovyanGFX für ESP32) sind generischen Lösungen (TFT\_eSPI) oft überlegen.
\item \textbf{Mehrere Tage für Display-Integration einplanen:} Die Kombination aus Hardware-Treiber, UI-Framework und Custom Fonts ist komplex und fehleranfällig.
\item \textbf{Schrittweise vorgehen:} Erst Display-Initialisierung, dann einfache Grafik, dann LVGL, dann Custom Fonts -- nicht alles gleichzeitig.
\end{enumerate}

%=============================================================================
\subsection{Lesson 2: Sensor-Anordnung im Gehäuse -- Airflow und Thermik}
\label{sec:lesson_gehaeuse}
%=============================================================================

\subsubsection{Ausgangssituation}

Das Gehäuse sollte alle Komponenten kompakt unterbringen und gleichzeitig präzise Messungen ermöglichen. Schnell zeigte sich, dass die räumliche Anordnung der Sensoren einen erheblichen Einfluss auf die Messgenauigkeit hat.

\subsubsection{Probleme und Lösungsweg}

\paragraph{Problem 1: Eigenerwärmung des MH-Z19C}
Der NDIR-CO\textsubscript{2}-Sensor MH-Z19C besitzt eine interne IR-Lampe und erwärmt sich im Betrieb um ca.\ 5--8\,°C über Umgebungstemperatur. In der ersten Gehäuseversion war der AHT20 (Temperatur/Feuchte) direkt neben dem MH-Z19C platziert. Die gemessene Temperatur lag dadurch konstant 3--4\,°C über der tatsächlichen Raumtemperatur.

\paragraph{Lösung: Thermische Zonierung}
Das Gehäuse wurde in zwei thermische Zonen aufgeteilt:
\begin{itemize}[nosep]
\item \textbf{Warme Zone (oben):} MH-Z19C, ESP32-S3, Display (alle wärmeabgebenden Komponenten)
\item \textbf{Kalte Zone (unten):} AHT20, SGP40, PMS5003, LD2410C (temperatursensitive Sensoren)
\end{itemize}

Eine Trennwand mit Lüftungsschlitzen sorgt für Konvektion, ohne direkten Wärmeübertrag.

\paragraph{Problem 2: Airflow für Feinstaubsensor}
Der PMS5003 benötigt aktive Luftzirkulation durch seinen internen Lüfter. In der geschlossenen ersten Version war die Luftzufuhr zu gering, was zu instabilen PM-Werten führte.

\paragraph{Lösung: Strömungsoptimierte Öffnungen}
Das finale PETG-Gehäuse hat:
\begin{itemize}[nosep]
\item Große Einlassöffnung (20$\times$30\,mm) an der Unterseite für PMS5003-Lüfter
\item Seitliche Schlitze (8$\times$40\,mm) für passive Konvektion
\item Abluftöffnung (15$\times$25\,mm) oben für Warmluftabfuhr
\end{itemize}

\paragraph{Problem 3: Radar-Abstrahlcharakteristik}
Der LD2410C hat einen horizontalen Erfassungswinkel von $\pm$60°. In der ersten Einbauposition (horizontal, parallel zur Gehäusewand) war die Reichweite auf unter 1\,m begrenzt. Durch 15°-Neigung nach außen wurde die spezifizierte Reichweite von 2\,m erreicht.

\subsubsection{Erkenntnisse}

\begin{enumerate}[nosep]
\item \textbf{Wärmeabgabe der Sensoren recherchieren:} Datenblätter geben oft die Verlustleistung an. Der MH-Z19C dissipiert ca.\ 200\,mW kontinuierlich.
\item \textbf{Prototyp ohne Gehäuse testen:} Zuerst die Sensoren im offenen Aufbau validieren, um Eigenerwärmungseffekte zu identifizieren.
\item \textbf{3D-Druck ermöglicht schnelle Iteration:} Das PETG-Gehäuse wurde dreimal überarbeitet. Der 3D-Druck war kostenlos verfügbar, was schnelle Iterationen ermöglichte.
\item \textbf{CAD-Simulation nutzen:} Strömungssimulation in Fusion~360 half bei der Optimierung der Luftöffnungen.
\end{enumerate}

%=============================================================================
\subsection{Lesson 3: Spannungspegel und Pegelwandlung -- Die 5V-Falle}
\label{sec:lesson_levelshift}
%=============================================================================

\subsubsection{Ausgangssituation}

Der ESP32-S3 arbeitet mit 3,3\,V-Logikpegeln. Mehrere Sensoren (MH-Z19C, PMS5003, LD2410C) werden jedoch mit 5\,V betrieben. Die Annahme, dass 5\,V-Sensoren auch 3,3\,V-kompatible Signale ausgeben, erwies sich als gefährlich falsch.

\subsubsection{Probleme und Lösungsweg}

\paragraph{Kritisches Problem: LD2410C TX-Pegel}
Der LD2410C-Radar arbeitet mit 5\,V-Logik auf seiner TX-Leitung. Ein erster Test ohne Level Shifter führte zu instabilem Verhalten des ESP32 -- Abstürze, spontane Resets, fehlerhafte Werte. Eine Messung mit dem Oszilloskop zeigte TX-Signale mit 5\,V-Pegel, die direkt auf GPIO\,7 gingen.

Nach Recherche in Espressif-Foren und Datenblättern stellte sich heraus: Der ESP32-S3 ist \textbf{nicht} 5\,V-tolerant. Maximale Eingangsspannung ist VDD+0,3\,V, also ca.\ 3,6\,V. Überspannung führt zunächst zu fehlerhaftem Verhalten, langfristig zur Zerstörung des GPIO-Pads oder des gesamten Chips.

\paragraph{Lösung: Bidirektionaler Level Shifter}
Ein TXS0108E-Modul wurde zwischen LD2410C und ESP32 geschaltet:
\begin{itemize}[nosep]
\item VA (Low Voltage): 3,3\,V vom ESP32
\item VB (High Voltage): 5\,V vom LD2410C
\item Automatische Richtungserkennung für TX/RX
\end{itemize}

\paragraph{Weitere 5V-Komponenten}
MH-Z19C und PMS5003 arbeiten ebenfalls mit 5\,V, geben aber glücklicherweise 3,3\,V-kompatible TX-Signale aus (Open-Drain mit Pull-up auf der Sensorseite). Dies wurde erst nach dem LD2410C-Problem systematisch überprüft.

\subsubsection{Erkenntnisse}

\begin{enumerate}[nosep]
\item \textbf{Datenblätter \textit{vor} dem Anschließen lesen:} Die kritische Information stand im LD2410C-Datenblatt, wurde aber zunächst übersehen.
\item \textbf{Oszilloskop verwenden:} Bei UART-Problemen immer die tatsächlichen Spannungspegel messen, nicht vermuten.
\item \textbf{Im Zweifelsfall Level Shifter einbauen:} Die 1,50\,\euro{} für ein TXS0108E-Modul sind billiger als ein neuer ESP32.
\item \textbf{Eingangsspannung der GPIOs kennen:} Nicht alle Mikrocontroller sind 5V-tolerant. Arduino Uno ja, ESP32 nein, STM32 teilweise.
\end{enumerate}

%=============================================================================
\section{Ausblick}
%=============================================================================

Das Projekt hat eine solide Basis für zukünftige Erweiterungen geschaffen:

\begin{itemize}[nosep]
\item \textbf{Cloud-Anbindung:} Die WebApp könnte um ein Backend erweitert werden, das Langzeit-Datenlogging und Push-Benachrichtigungen ermöglicht.
\item \textbf{Weitere Sensoren:} Der ESP32-S3 bietet noch freie GPIOs für zusätzliche Sensoren (z.\,B.\ Ozon, NO\textsubscript{2}).
\item \textbf{Batteriebetrieb:} Mit einem LiPo-Akku und optimiertem Deep Sleep könnte eine portable Version entstehen.
\item \textbf{Individuelles PCB:} Der aktuelle Aufbau auf Lochrasterplatine könnte durch eine professionelle Platine ersetzt werden.
\end{itemize}

Das V-Modell hat sich als geeignete Methodik für dieses Embedded-Projekt bewährt. Die systematische Spezifikation und Verifikation hat dazu beigetragen, dass alle Anforderungen erfüllt und dokumentiert wurden.
